\section{Theoretical Foundations of Semantic Routing}
\label{sec:theory}

Section~\ref{sec:method} defines a simple routing rule from a stage response
$r(t)$ and final contrast $d$ to evidence, contradiction, and fragility.
We first ask whether this routing is arbitrary.

To formalize this routing, we establish three operational axioms that reflect practical explanation goals. First, \emph{conservation} ensures that no response magnitude is lost or invented. Second, \emph{endpoint gating} isolates sensitivity on pairs whose final contrast is negligible under the chosen threshold. Third, \emph{directional support} distinguishes intermediate responses that move toward the model's final decision from those that oppose it. 

\begin{theorem}[Unique hard-gated semantic routing]
\label{thm:canonical}
Fix $a\in\{0,1\}$, endpoint orientation $s\in\{-1,1\}$ on the active branch, and response $r\in\mathbb{R}$. The unique nonnegative triple satisfying \textbf{conservation}, $e+c+f=|r|$; \textbf{endpoint gating}, $f=0$ for $a=1$ and $e=c=0$ for $a=0$; and \textbf{directional support}, $c=0$ when $sr\geq0$ and $e=0$ when $sr\leq0$, is
\[
(e,c,f)=\bigl(a(sr)^{+},\;a(sr)^{-},\;(1-a)|r|\bigr).
\]
\end{theorem}

Theorem~\ref{thm:canonical} removes tunable mixtures once the endpoint gate and orientation are specified. It establishes uniqueness within the stated endpoint-relative semantics, not uniqueness of hard endpoint gating itself. Under these operational axioms, the routing is uniquely determined. \emph{See Appendix~\ref{app:theory} for the proof and the positive--negative (Jordan) decomposition interpretation.}

Theorem~\ref{thm:canonical} shows that the routing is unique once the endpoint-relative semantics are fixed. We next ask whether retaining the three routed components provides information beyond ordinary magnitude.
% \paragraph{Magnitude non-identifiability.}
\begin{theorem}[\decaf strictly refines response magnitude]
\label{thm:nonident}
Ordinary response magnitude is a deterministic function of the \decaf
profile,
\[
\mathrm{Abs}=E+C+F,
\]
so every decision rule based on $\mathrm{Abs}$ can also be implemented from
$(E,C,F)$. The reverse recovery fails in general: for every $m>0$, the
distinct profiles $(m,0,0)$, $(0,m,0)$, and $(0,0,m)$ all produce
$\mathrm{Abs}=m$. Hence the refinement is strict for any decision problem
that distinguishes two such profiles with positive probability.
\end{theorem}
Thus \decaf is a lossless refinement of ordinary magnitude: any analysis
based on $\mathrm{Abs}$ can be reproduced from $(E,C,F)$, while the reverse
recovery is impossible in general. The refinement is strict whenever a
decision problem distinguishes response profiles that share the same
magnitude. For example, if evidence, contradiction, and fragility are equally
likely and all produce magnitude $m$, every magnitude-only classifier receives
the same observation and cannot exceed $1/3$ accuracy, whereas the full
profile distinguishes the three roles. Section~\ref{sec:imagenet9} tests this consequence empirically after explicitly
matching ordinary response magnitude.
See Appendix~\ref{app:theory} for the Bayes limit, strict information refinement, and unequal priors.

% \paragraph{Preservation, attenuation, and inversion.}
Theorem~\ref{thm:nonident} shows that magnitude can collapse distinct response roles. A practically important case is the difference between an effect that disappears and one that reverses. Both can reduce aligned evidence, but only reversal should create contradiction.

Suppose the final contrast is active with magnitude $m>0$.
 Let $\eta\in[0,1]$ denote the probability that this effect is altered by context. With probability $1-\eta$, the response remains aligned with magnitude $m$; with probability $\eta$, it is either suppressed to zero (attenuation) or reversed to $-m$ (inversion). We consider three simple response regimes: \emph{preservation}, where the effect keeps its direction; \emph{attenuation}, where it weakens toward zero; and \emph{inversion}, where it reverses direction. 

\begin{proposition}[Contradiction separates attenuation from inversion]
Under equal aligned and opposed magnitudes,
\[
{\footnotesize
\begin{array}{lccc}
\toprule
\text{Regime} & E & C & \mathrm{Abs} \\
\midrule
\text{Preservation}       & m           & 0      & m \\
\text{Attenuation to zero}& (1-\eta)m  & 0      & (1-\eta)m \\
\text{Inversion}          & (1-\eta)m  & \eta m & m \\
\bottomrule
\end{array}
}
\]
In the inversion regime, $C/(E+C)=\eta$.
\label{prop:calibration}
\end{proposition}
Attenuation and inversion can remove the same amount of aligned evidence, but only inversion creates contradiction. Thus C distinguishes loss of an effect from reversal of that effect. The controlled experiment in Section~\ref{sec:controlled}.
\emph{See Appendix~\ref{app:theory} for the proof, unequal magnitudes, and continuous mixtures.}

When contradiction and fragility vanish, conservation gives $\text{Abs}=E$. Thus DECAF reduces exactly to ordinary magnitude when the response is already fully aligned; it introduces additional distinctions only when the observed behavior contains them.
